% TOP_PROBLEM: {"id": 2305071, "title": "Research Problems in Function Theory — Problem 5.71", "category_id": 8, "category": {"id": 8, "name": "analysis", "display_name": "Analysis", "description": "Limits, continuity, calculus, and function theory.", "slug": "analysis", "order_index": 8, "created_at": "2026-07-31T15:26:25.671Z"}, "status": "open", "problem_number": "AMR-022-5071", "set_id": 13}
% FIRST_POSTED: 2026-10-02
% ENTRY_KIND: statement_only
% UnsolvedMath ID 2305071; source catalogue code AMR-022-5071
% !TeX program = lualatex
% Definition and sources only; this file is not an LLM solution attempt.
\documentclass[11pt,a4paper]{article}
\usepackage[margin=25mm]{geometry}
\usepackage{fontspec}
\setmainfont{lmroman10-regular.otf}[BoldFont=lmroman10-bold.otf,
 ItalicFont=lmroman10-italic.otf,BoldItalicFont=lmroman10-bolditalic.otf]
\usepackage{amsmath,amssymb,mathtools}
\usepackage{unicode-math}
\setmathfont{latinmodern-math.otf}
\AtBeginDocument{\let\setminus\smallsetminus}
\usepackage{xurl,hyperref}
\hypersetup{colorlinks=true,urlcolor=blue}
\setcounter{secnumdepth}{0}
\setlength{\parindent}{0pt}
\setlength{\parskip}{5pt}
\setlength{\emergencystretch}{3em}
\begin{document}
\section*{Research Problems in Function Theory — Problem 5.71}\label{2305071}
{\small UnsolvedMath ID 2305071 · AMR-022-5071 · Analysis · AMR Open Problem Lists}
\subsection{Definitions and mathematical statement}
Let $f(z)=\sum_{k\ge1}a_kz^{n_k}$ be holomorphic in the unit disc, where the nonzero exponents are positive integers satisfying $n_{k+1}/n_k\ge q>1$ for a fixed $q$. Define
\[
T(r,f)=\frac1{2\pi}\int_0^{2\pi}\log^+|f(re^{i\theta})|\,d\theta,
\quad
m(r,w;f)=\frac1{2\pi}\int_0^{2\pi}
\log^+\frac1{|f(re^{i\theta})-w|}\,d\theta.
\]
Assume $T(r,f)\to\infty$ as $r\uparrow1$. The deficiency of a finite value $w\in\mathbb C$ is
$\delta(w,f)=\liminf_{r\uparrow1}m(r,w;f)/T(r,f)$.
Must $\delta(w,f)=0$ for every finite $w$?
The question concerns the relative rarity measured by integrated logarithmic means; merely proving that each value is assumed infinitely often is weaker. The gap condition holds with one ratio $q$ for the complete exponent sequence. The coefficients may tend to zero, provided their combined growth still makes $T$ unbounded.
Primary work of Tsuboi gives additional sufficient coefficient-growth conditions and recalls Murai's result when the coefficients do not tend to zero. Those special cases do not remove the general unbounded-characteristic hypothesis posed here.
\subsection{Short English statement}
Does every disc power series with Hadamard gaps and unbounded characteristic have no finite deficient value?
\subsection{Sources}
\begin{itemize}
\item[S1] ULAM.AI and the UnsolvedMath Contributors, supplied problems.json, record 2305071 (AMR-022-5071), AMR Open Problem Lists. Catalogue identity and source transcription; CC BY 4.0.
\url{https://huggingface.co/datasets/ulamai/UnsolvedMath}.
\item[S2] Hayman - Research Problems in Function Theory (2018), Problem 5.71. Original source indexed by AMR-022-5071.
\url{https://arxiv.org/abs/1809.07200}.
\item[S3] N. Tsuboi, The defects of power series in the unit disk with Hadamard gaps, introduction and sufficient conditions.
\url{https://www.cajpn.org/complex/refs/thesis/16D-Tsuboi.pdf}.
\end{itemize}
\end{document}
